Quantum algorithms offer the potential to address combinatorial optimisation problems that are intractable for classical computers. Among these, the graph similarity problem remains a fundamental challenge with applications across network science, logistics, and bioinformatics. This thesis investigates the performance of the non-variational quantum walk-based optimisation algorithm (NVQWOA) on this problem.

Through state-vector simulations, this work characterises the behaviour of the NVQWOA for multiple problem instances, sizes, and algorithm iterations. The algorithm consistently amplifies the probability of optimal and near-optimal solutions relative to the uniform initial state and outperforms Grover's search algorithm on an equivalent number of iterations, demonstrating that constructive interference concentrates amplitude in regions of high solution quality. Analyses based on Hamming and subshell distance reveal that amplification strength depends on structural proximity to the optimum, reflecting the locality of the quantum walk on the permutation graph. Techniques for numerically optimising the hyperparameters are investigated. Additional investigation of quality gaps and variance trends shows that the solution space structure enhances the effectiveness of amplitude redistribution and satisfies key conditions for efficient optimisation.

These results establish NVQWOA as a promising approach to solving the graph similarity problem. Its sensitivity to solution structure and minimal optimisation requirements suggest potential advantages for near-term quantum devices.